\title{AI-Native Firms}
\begin{document}
\maketitle

\begin{abstract}
Hyunjin Kim Rembrand Koning Working Paper 26-090
\end{abstract}

\section{Recovered Text}
Working Paper 26-090
AI-Native Firms
Hyunjin Kim
Rembrand Koning
Working Paper 26-090
Copyright © 2026 by Hyunjin Kim and Rembrand Koning.
Working papers are in draft form. This working paper is distributed for purposes of comment and discussion only. It may
not be reproduced without permission of the copyright holder. Copies of working papers are available from the author.
Funding for this research was provided in part by Harvard Business School.
AI-Native Firms
Hyunjin Kim
INSEAD
Rembrand Koning
Harvard Business School
AI-Native Firms∗
Hyunjin Kim
INSEAD
Rembrand Koning
Harvard Business School
June 9, 2026
We study how firms built around AI capabilities—“AI-native” firms—are organized. Drawing
on Y Combinator batches W20–F24 and U.S. venture-backed startups whose first financing
closed between 2020 and 2024, we classify each firm’s AI-native status and link it to workforce
microdata on team size, function, seniority, and hierarchy. Relative to non-AI startups in the
same industry-cohort, AI-native firms are 25\% smaller. Their share of engineers is 13\% greater,
and the shares of entry-level workers and managers are each roughly 15\% lower. Their hierarchies
are half a seniority level flatter—yet valuations are comparable, suggesting higher value created
per employee. We argue these patterns reflect two channels: a process channel, in which AI
changes how people work inside the firm, and a product channel, in which AI capabilities are
built into what the firm sells. Using text from product descriptions and job postings, we find
that embedding AI into the product, beyond layering on AI tools into existing workflows, is a
primary way startups are scaling “knowledge work” without large teams of knowledge workers.
∗Kim: h@insead.edu. Koning: rem@hbs.edu. The authors gratefully acknowledge financial support from the AI
Institute at Harvard Business School and INSEAD. The authors thank Samay Gupta and Simonne Pinto for research
assistance. We thank Ajay Agrawal, Arnaldo Camuffo, Lauren Thomas, Todd Lensman, and Abhishek Nagaraj along
with seminar participants at SMS, the HBS EM working lunch, and the INSEAD brownbag for their especially helpful
feedback.
1
1
Introduction
Technology reshapes organizational form. Railroads and the telegraph allowed firms to serve na-
tional markets, increasing the returns to coordinating across distance and creating the conditions
for large managerial hierarchies (Chandler, 1993). Electricity altered how work could be organized
on the factory floor and spurred the rise of scientific management (David, 1990).
The internet
lowered communication and coordination costs, supporting more platform-based, networked, and
distributed forms of organization (Shapiro and Varian, 1999). Across technologies—from lithog-
raphy to Linux—product and organizational architectures often co-evolve (Henderson and Clark,
1990; Colfer and Baldwin, 2016).
How will AI reshape firms? Recent research has largely focused on a process channel, empha-
sizing how AI tools augment or automate work inside the firm through agentic coding, AI-assisted
customer service, automated sales operations, and many other internal tasks (Peng et al., 2023;
Brynjolfsson, Li, and Raymond, 2023; Dell’Acqua et al., 2023; Csaszar, Ketkar, and Kim, 2024;
Otis et al., 2025; Sarkar, 2026; Kim, Kim, and Koning, 2026). Yet AI can be used for more than
internal work. AI-native firms embed AI into their products, allowing customers to generate slide
decks with AI, build software, and even run their accounts receivable. This product channel shifts
productive capability from the internal organization into the product, moving knowledge work that
would have required large teams to be performed directly in the product itself through imported
model capabilities. Both channels suggest that “AI-native” firms—those that are built around AI-
enabled products and processes—may be organized differently than comparable non-AI firms.
In this paper, we highlight this product channel and explore whether AI-native firms are, in
fact, organized differently. We focus on startups because new firms are less constrained by legacy
structures, making them a likely setting in which new organizational forms first become visible.
Measuring whether a firm is AI-native presents a challenge because, as our two channels illustrate,
there are many ways companies can be built around AI. Here, we leverage the fact that Y Combi-
nator (YC), a leading accelerator for early-stage companies, curates and encourages startups to tag
themselves with public labels to help investors, journalists, and customers discover new ventures
they might be interested in. We leverage the fact that startups can add an “AI” tag to their public
profile, enabling us to identify startups that claim that AI is central to their firm, whether that is
embedding it into their products, using it as a core part of their production process, or something
else entirely. To test if our YC results generalize beyond this extremely selected set of firms, we also
2
use text describing each YC startup to train a model that predicts any startups’s AI status, allowing
us to extend our analysis to the broader universe of venture-backed U.S. startups in PitchBook. We
then link both samples to Revelio Labs for detailed workforce data on team size, hierarchy, seniority,
and functional composition.
We structure our analysis into two parts. First, we document a robust set of organizational
differences between AI-tagged and non-AI-tagged startups. We find that AI-tagged startups have
approximately 25\% fewer employees than non-AI firms in the same industry-cohort cell.
Their
workforces are also more technical: the share of engineers is 5 percentage points higher, offset by
lower shares in sales, finance, operations, and administration. They skew senior, with the share of
entry-level workers just under 15\% lower and the share of senior workers about 20\% higher. Their
hierarchies are flatter: half a level less depth and 15\% fewer managers, even controlling for firm size.
These patterns are consistent across our YC and PitchBook samples.
These smaller teams raise similar amounts of funding and hit valuations on par with non-AI-
native firms, suggesting that they are not lower quality entrants. As a consequence of being smaller,
AI-tagged startups raise roughly 20\% more capital per employee and carry higher valuations per
employee. These smaller teams also draw from a narrower labor pool: they are geographically con-
centrated in Silicon Valley, and employ workforces that are more male, more likely to hold advanced
degrees, and drawn from more prestigious employers and institutions. Our findings are consistent
with AI complementing technical expertise at the firm level, with small teams of skilled founders
building ventures that historically might have required much larger hierarchical organizations.
The second part of our analysis explores the role of the product and process channels in explaining
these patterns. Focusing on our YC sample, for which we have richer text data, we use an LLM
to inductively classify founders’ own descriptions of their businesses and ask how AI is embedded
into the product. 43\% of AI-tagged YC startups use AI to fully automate tasks workers used to do,
24\% build AI tools designed to augment existing workers at the firms they sell to, and 15\% build AI
infrastructure, providing a capability that another developer integrates into their own AI product.
To measure the process channel, we code whether each firm’s job postings name specific worker-
facing AI tools like ChatGPT, GitHub Copilot, or Cursor. AI-native startups are 2.6 times as likely
as non-AI startups to name such tools in their postings, but this measure does not predict smaller
firm sizes in the log, Poisson, or hierarchy specifications. Even when controlling for this process
channel, we find that embedding AI into a startup’s products is associated with less headcount and
flatter structures.
3
If embedding AI into the product moves productive capability from the firm’s workers into
the product itself, then the firms that should shrink most are those that, without AI, would have
scaled by hiring people to deliver the underlying work—therapy, tutoring, telemedicine, and similar
service businesses where small teams can now build AI capable of doing the work that had required
building larger workforces. To test this idea, we test for heterogeneity in our estimates by whether
the startup is classified as a “services” business. The difference between AI and non-AI firms is
significantly larger for these services firms: AI startups are roughly 70\% smaller than non-AI peers
and have hierarchies nearly a full level flatter.
Our paper makes three contributions. First, we identify a product channel largely overlooked
in recent work on AI’s productivity impacts. The dominant research paradigm today focuses on
the process channel: how AI tools change the firm’s internal workflows (Brynjolfsson, Li, and
Raymond, 2023; Dell’Acqua et al., 2023; Kim et al., 2024). But a key way in which firms leverage
AI is to embed it directly into what they sell, so that model capabilities perform, augment, or
coordinate knowledge work through the product itself. This matters because it changes where the
firm’s productive capability resides. A long line of work models the firm as a system for processing
information and knowledge (Galbraith, 1973; March and Simon, 1958; Simon, 1947; Garicano, 2000),
and firm performance has been shown to depend on how organizations access and act on decision-
relevant information (Kim, 2025). Our findings suggest that AI enables firms to shift some of this
processing from the organization into the product itself, allowing the firm to deliver knowledge-work
capabilities without building them internally through people and hierarchy (Ide and Talamás, 2025).
If this pattern holds, then building, importing, and orchestrating model capabilities may matter as
much for competitive advantage as building human ones (Nelson, 1985; Teece, Pisano, and Shuen,
1997; Helfat and Peteraf, 2003).
Second, we show that these AI-tagged firms employ smaller teams of more talented and technical
workers. These workers are especially likely to be graduates from elite institutions, concentrated in
Silicon Valley, and male. While at the task-level AI often helps less experienced workers perform
better (see Brynjolfsson, Li, and Raymond, 2023; Noy and Zhang, 2023; Csaszar, Ketkar, and Kim,
2024, ; cf. Otis et al. 2025), at the firm level we see a shift broadly in the opposite direction. Existing
work shows generative AI changes labor demand when adopted by existing firms (e.g., Gupta et al.,
2025; Brynjolfsson, Chandar, and Chen, 2025); we show that AI also changes the kind of firms that
are created. While this may have implications for career opportunities and labor market outcomes
(Beane, 2024; Kang and Kim, 2025), we urge caution in moving from firm-level to market-level
4
estimates. While AI-tagged startups employ fewer workers, they appear much less costly to launch
and scale. If AI leads to meaningfully more and faster growing startups (Gupta et al., 2025), then
while each firm may be smaller the overall demand for labor could well rise.
Third, our finding on the product and process channels is consistent with emerging evidence on
the “mapping problem” in AI adoption. Kim, Kim, and Koning (2026)’s field experiment shows that
the binding constraint to AI adoption is not access to the technology but the ability to search across
a firm’s own processes and discover where AI can create value. Firms that solve this problem grow
faster while requiring less external capital and no additional labor. Our null process-channel result,
though a coarse and noisy proxy, is consistent with this view: equipping workers with ChatGPT,
Copilot, or Cursor does not, on its own, predict smaller firms. Capturing real value from AI often
requires re-engineering processes around it—and, as our results suggest, re-engineering the product
itself so that AI does the work directly. The AI-native ventures in our data appear to have solved
this problem by building the firm around AI from the start. More broadly, our findings suggest
that when firms are fully built around AI, the technology does more than lower the cost of entry
(Ewens, Nanda, and Rhodes-Kropf, 2018; Reimers and Waldfogel, 2026), it changes how ventures
are organized and scale.
The remainder of the paper proceeds as follows. Section 2 elaborates on what makes a firm
AI-native, distinguishing the process and product channels. Section 3 describes our two comple-
mentary datasets—Y Combinator startups and the broader universe of U.S. venture-backed firms
from PitchBook—along with workforce data from Revelio Labs. Section 4 presents results on firm
size, functional composition, seniority, hierarchy, and capital efficiency. Section 5 then explores the
role of the product and process channels in explaining these results. Finally, section 6 discusses
implications and limitations.
2
What Makes a Firm AI-Native?
The term “AI-native” firms has emerged in industry discussions to refer broadly to firms built
around AI capabilities. Industry reports suggest that AI startups may be growing with smaller
teams, but provide widely varying estimates and limited insight into how organizational form is
changing (Tables A1 and A2). One reason may be that the term “AI-native” is used to describe
several distinct ways in which firms rely on AI. In some accounts, AI-native firms are those that use
AI as an internal production technology or core operating model, changing how employees code, sell,
5
design, or coordinate work. We refer to this as the process channel. In other accounts, AI-native
firms are those that produce products and services that would not have been possible without AI,
ranging from foundation model providers to companies that embed AI directly into what they sell.
We refer to this as the product channel.
The process channel has been the focus of most existing research, which has largely focused on
how AI can either augment or automate tasks within firms. The majority of recent research on
the impact of generative AI on workers and firms has focused on this channel, exploring how AI
changes the productivity of workers on individual tasks like coding, customer service, and strategy
development (Peng et al., 2023; Brynjolfsson, Li, and Raymond, 2023; Csaszar, Ketkar, and Kim,
2024).
This work has sparked speculation about a process channel by which organizations can
succeed with smaller teams equipped with AI tools that enable workers to match the productivity
of much larger teams. Building on these studies exploring the impact of AI on individual tasks,
recent experimental evidence suggests that AI can have firm-level impacts on performance, growth,
and productivity, but only when firms know how to map it into production (Kim, Kim, and Koning,
2026).
Alongside such internal use, firms are also increasingly embedding AI into their products. The
organizational implications of this product channel are not immediate: adding AI features to a prod-
uct may leave the firm’s structure largely unchanged, much as shifting from one web development
language to another is unlikely to alter the firm’s underlying structure or production function. The
implications are different, however, if AI allows the product itself to perform knowledge work that
would otherwise require human teams. In this case, capabilities that once had to be built inside the
organization can be embedded directly into what the firm sells.
A useful way to understand this is to start from the view of the firm as an information-processing
system. Firms develop routines and hierarchy to help workers process information and coordinate
effectively (Simon, 1947; March and Simon, 1958; Galbraith, 1973; Garicano, 2000). These struc-
tures make firm growth possible, but they also tie growth to a human constraint: knowledge must be
allocated and processed through people and hierarchy (Garicano and Rossi-Hansberg, 2015), deci-
sions require judgment under uncertainty (Agrawal, Gans, and Goldfarb, 2022), and organizational
performance depends on the scarce attention and coordination needed to translate decisions into
action (Dessein, Lo, and Minami, 2022; Kim, 2025). Growth therefore typically requires expanding
the human organization through which knowledge work is performed.
When AI is embedded in the product, some of this knowledge work can be performed directly
6
through what the firm sells. Producing personalized recommendations for different types of cus-
tomers no longer requires a team to segment users, tailor messages, and coordinate around the task,
since AI can perform this work inside the product. Providing credit score estimates from within
the product does not require a division of workers producing such predictions. AI thus enables
capabilities that once required repeated human workflows inside the organization to move into the
product, potentially changing how firms scale and compete (Nelson, 1985; Teece, Pisano, and Shuen,
1997; Helfat and Peteraf, 2003).
This argument builds on Ide and Talamás (2025)’s formal model of AI’s impact on knowledge
hierarchies within firms. Building on hierarchical models of knowledge work (Garicano, 2000), they
model AI as making noncodifiable expertise scale through compute rather than relying only on
scarce human time and knowledge. They distinguish between nonautonomous AI, which advises
human workers, and autonomous AI, which can pursue production opportunities independently.
In their model, autonomous AI changes production because routine work can be performed by AI
agents rather than human workers, reallocating humans across production and problem-solving roles
and giving high-knowledge workers greater leverage.
Building on this insight, we focus on what happens when autonomous AI is embedded in the
product rather than deployed inside the firm’s production process. Scaling outputs through knowl-
edge work traditionally requires organizing more production internally: employees scope requests,
coordinate on work, and deliver results. Autonomous AI can make each workflow less labor-intensive
and shift employees toward higher-level tasks, but producing more outputs still requires the firm to
manage more work. In contrast, when AI is embedded in a product, customers can generate out-
puts directly through the product. Users specify the task, interact with the AI, evaluate and revise
results, and iterate directly in the interface. While the firm still invests in building the product,
output growth can be scaled through compute rather than by expanding internal teams.
This shift may have implications for how firms are organized. First, it may lower the need for
human labor per unit of output, because serving additional customers does not require a propor-
tional increase in internal production work. Second, it may contribute to flatter structures, because
customer requests can be handled through the product rather than routed through a knowledge hier-
archy. Third, it may increase technical talent density, as larger shares of employees focus on building
and improving the product rather than carrying out routine production tasks. More broadly, pro-
ductizing capabilities that would have otherwise been delivered through internal workflows can
change how much labor firms need, as well as what kind of people and organizational design they
7
need to scale.
Recent cases illustrate how building around AI may shift how firms are organized. Gamma,
an AI presentation startup, builds on image and text foundation models to generate complete
slide decks for its users (Koning, Keller, and Kim, 2025). While Gamma’s engineers use AI tools
extensively, including tools such as Cursor for software development, the case illustrates how the
largest organizational implication comes from embedding AI into the product itself. One way to
view Gamma’s product is as the deployment of millions of autonomous or semi-autonomous AI
agents, each pursuing a production opportunity on behalf of a user: generating layouts, writing
copy, selecting images, and localizing content across dozens of languages. In contrast, a traditional
presentation service company, even one using AI internally, would still scale by turning each deck
request into an internal workflow: scoping the request, coordinating design and copy work, and
reviewing revisions before delivering the final deck. Gamma instead turns the request into a product
interaction. The firm still needs employees to build and improve the product, but it does not need
to organize a human workflow around each deck. As a result, a team of roughly 30 employees could
serve millions of users and reach \$50 million in annual revenue in just two years.
FazeShift shows how the same product channel logic operates in the enterprise sector (Koning,
Kim, and Keller, 2025). The company builds AI agents that automate the full accounts receivable
lifecycle, including invoicing, collections, and cash application—work that previously required ded-
icated AR teams at each customer. With a team of just ten, FazeShift landed dozens of enterprise
customers by building agents that match payments to invoices, send context-aware collection mes-
sages, and resolve discrepancies that human analysts used to handle manually. A traditional AR
SaaS vendor would need a much larger team both to build out the firm-specific exception handling
for the AR product and to staff the operations to resolve the exceptions that FazeShift’s AI agents
now handle automatically.
Taken together, these arguments suggest AI-native firms should look different. On the process
side, we should expect such firms to be more productive, with smaller teams able to do the same
amount of work—although if AI dramatically lowers the cost of what the firm sells, a Jevons-style
expansion in demand could lead firms to hire more rather than fewer workers. On the product side,
as firms move AI capabilities into what they sell, they may require different sorts of talent: workers
with the technical expertise to embed AI into a product and to complement, rather than be replaced
by, those capabilities. Both channels suggest that AI-native firms should be organized differently
from their non-AI peers, as the technology changes both how the firm works internally and what
8
the firm needs to build. In the next section, we turn to the data to ask whether AI-native firms do
in fact look different, and how each channel—process and product—account for the differences we
find.
3
Data and Research Design
We construct two complementary datasets to examine this question. Our primary dataset focuses
on startups that attend YC, which provides us with clean measurement of AI status and startup
age. Our secondary dataset is PitchBook, which allows us to extend our analysis to the broader
universe of venture-backed startups in the U.S. to test the external validity and generalizability of
our core YC estimates.
3.1
Y Combinator Sample
We construct a dataset comprising all YC startups that participated in batches from Winter 2020
to Fall 2024, collected through an open-source public API.1 Since YC publishes complete rosters for
each cohort, it offers a well-defined sample through which hundreds of young but promising compa-
nies pass each year. The 2020–2024 window spans GPT-3’s public release in 2020 and ChatGPT’s
breakout in 2022—the period when AI funding and startup formation surged. Our initial sample
contains 2,891 firms across 11 YC cohorts (W20, S20, W21, S21, W22, S22, W23, S23, W24, S24,
F24).
3.1.1
AI Classification
YC encourages startups to tag themselves with descriptive labels to facilitate discovery by journal-
ists, potential users, and investors. A startup like FazeShift, an “AI agent for Accounts Receivable,”
is marked with tags including “artificial-intelligence,” “fintech,” and “b2b.” We classify a startup
as an AI startup if its tags include “artificial-intelligence” or “ai.” Reflecting the general-purpose
nature of the technology, AI startups build around AI in varied ways: from autonomous agents to
AI driven drug discovery to AI-guided physical work. Self-identification lets us rely on the founders’
beliefs about whether AI is central to their firm, providing a broad measure of whether the firm is
built around AI or not.2 Especially given YC’s prominence and the reputational stakes involved,
1API available at https://github.com/yc-oss/api.
2The idea of building firms around AI could allow smaller firms to scale predates the launch of ChatGPT in late
2022. In 2017, five years before ChatGPT, a cofounder of InsureTech startup Lemonade urged founders to “don’t
9
there are strong incentives for these public statements to accurately reflect the firm’s underlying
use of the technology.3 Of the 2,891 YC startups in our sample, 990 (34.2\%) are tagged as AI
startups and 1,901 (65.8\%) as non-AI startups (see Figure 1 for the distribution across batches
in the PitchBook-matched regression sample and Appendix Figures A1–A2 for example YC pages
illustrating this tagging).
3.1.2
Linking to PitchBook
We merge our YC data with PitchBook, a private-market data intelligence platform, using nor-
malized company website URLs.
PitchBook provides comprehensive data on startup financing,
including deal-level funding amounts, valuations, and deal types.4 Of our 2,891 YC firms, we suc-
cessfully match 2,786 (96.4\%) to PitchBook records. The small number of unmatched firms are
primarily very early-stage companies without substantial financing activity or companies that have
since become inactive.
Appendix Table A3 consolidates the match rates for both the YC and
PitchBook sample construction.
3.1.3
Linking to Revelio Labs
To measure organizational structure and firm size, we link our YC-PitchBook data to Revelio
Labs, a workforce intelligence platform that aggregates information from public professional profiles
and job postings to create comprehensive employee records. We use LinkedIn profile URLs from
PitchBook to match companies to a pull of Revelio’s database covering records through January
2025, successfully linking 2,233 firms (77.2\% of the full YC sample). Match rates are nearly identical
for AI and non-AI startups (76.2\% vs. 77.8\%) and range from 69\% to 82\% across batches. Given
that these are knowledge-work startups whose employees overwhelmingly maintain LinkedIn profiles,
Revelio’s inability to link any employees to a firm strongly suggests the firm has effectively zero
employees as of January 2025, so we impute zero employees for firms with no Revelio-matched
workers.
As our intensive margin results show, our findings are robust to coding non-matched
hire more people, build smarter AI” and described his own firm as one designed around “fewer people and more
code” (Wininger, 2017). As we discuss below, Lemonade’s vision mirrors our product channel. For example,
embedding an ML system for predicting fraud into the product replaces the large teams that would otherwise
manually review cases. Moreover, while ChatGPT led to a surge of AI startups in 2023, the first AI coding assistant
is GitHub’s Copilot, which was launched in 2021. OpenAI’s GPT-3 was launched in 2020.
3Of course, self-disclosure risks strategic manipulation, similar to firms adding “.com” to their names during the
internet bubble (Cooper, Dimitrov, and Rau, 2001). However, if the AI label is only being used superficially and
symbolically, then while we might see exuberant or irrational shifts in valuation it seems unlikely that it would
result in changes in how the firm is organized as measured by a third party data source (LinkedIn).
4That said, valuation and funding amounts are often missing for early stage firms like those we analyze.
10
startups as having missing values instead. Revelio provides position-level data including job titles,
seniority levels, start and end dates, and functional classifications. Appendix Table A3 reports
per-batch and field-level coverage rates.
3.2
PitchBook Sample: Extending to the Broader Startup Ecosystem
While YC provides a measurement of a startup’s AI focus and Revelio rich organizational data, it
represents a highly selected sample of startups. To test whether patterns observed in YC generalize
to the broader startup ecosystem, we extend our analysis to the universe of early-stage venture-
backed startups using the PitchBook data.
3.2.1
Sample Construction
Appendix Table A3 (Panel B) reports the PitchBook side of our sample-construction funnel. We
begin with all U.S.-based startups in PitchBook with a first financing deal between 2020 and 2024
(inclusive). The 2020–2024 first-deal window aligns with the YC W20–F24 sample period and the
January 2025 Revelio employment snapshot, focusing on the post-GPT-3 era and ensuring every
firm in the sample had at least one financing event before our outcome is measured. Using first
deal date (rather than founding year) allows us to construct deal quarter fixed effects comparable
to YC’s quarterly batch indicator.
3.2.2
Predicting AI Status
Unlike YC, PitchBook does not provide a clean indicator of whether a startup team wants to publicly
disclose if it is an AI startup.5 To classify startups as AI or non-AI, we train a machine learning
model on the YC sample and then use the model to predict AI status for PitchBook companies
using text descriptions of each startup.
In particular, we use PitchBook’s text descriptions of each company to generate 3,072-dimensional
embeddings using OpenAI’s text-embedding-3-large model. We then train a Random Forest classi-
fier to predict YC’s AI tags from these embeddings. We use 5-fold cross-validation to tune hyper-
parameters. The model achieves solid if unspectacular out-of-sample performance: cross-validated
AUC of 0.81, and on a held-out test set, accuracy of 77.5\%, precision of 73.9\% (when we predict
AI > 0.5, 74\% are actually AI), and recall of 55.2\% (we identify 55\% of true AI companies). In
5PitchBook assigns companies to “verticals” including an “Artificial Intelligence \& Machine Learning” category, but
this classification appears noisy, substantially undercounting AI startups relative to our YC-based measure. It is
also unclear who applies this measure to the startup and what the incentives are for truthful reporting.
11
the paper’s regressions, we use the continuous predicted probability rather than a binary threshold
as it provides more information. For descriptive statistics splitting the sample into AI and non-AI
groups, we use a 0.5 threshold.
3.3
Variable Definitions
AI Status.
For the YC sample, AI status is a binary indicator equal to 1 if the startup’s YC
tags include “artificial-intelligence” or “ai.” For the PitchBook sample, AI status is the predicted
probability from our Random Forest classifier trained on YC data. Since both variables range from
zero to one, it is straightforward to compare effect sizes across the two.
Employee Counts.
We measure employee counts using Revelio Labs data as of January 2025,
which captures employees matched through public professional profiles. While our primary analyses
use the January 2025 measures, we also construct a quarterly panel of employment counts to shed
light on how firm size changes over time for AI and non-AI tagged firms.
Manager Share.
The share of employees with O*NET occupation codes beginning with “11-”
(Management Occupations), based on Revelio’s mapping of job titles to standardized occupational
categories.
Hierarchy Depth.
The count of distinct Revelio seniority levels (1–7) present among a firm’s
employees. To calculate a meaningful hierarchy measure, we require at least one employee in Revelio
data, yielding a sample size of 2,233.
Funding and Valuation Per Employee.
Total funding raised divided by employee count (fund-
ing per employee) and last known valuation divided by employee count (valuation per employee). As
mentioned, funding and valuation amounts are often missing. For per-employee financing analyses,
this yields 2,218 YC firms and 29,336 PitchBook firms with funding data, and 1,046 YC firms and
12,891 PitchBook firms with valuation data.
Seniority Shares.
The share of employees at each seniority level based on Revelio Labs’ 1–
7 seniority scale: Entry-level (1–2), Mid-level (3–5), and Senior (6–7). To calculate meaningful
seniority share measures, we require at least one employee in Revelio data, again yielding a sample
size of 2,233 firms.
12
Functional Shares.
The share of employees in each functional category based on Revelio Labs job
classifications: Engineering, Science, Sales, Marketing, Operations, Finance, and Administrative.
Industry Group.
PitchBook assigns each company to one of 41 industry groups based on its
primary business activity; examples include Application Software, Information Services, Healthcare
Services, and Financial Software. We use PitchBook’s industry classification rather than standard
SIC codes because SIC coverage is extremely sparse for early-stage startups. In our YC sample,
fewer than 1\% of firms (24 of 2,786) have SIC codes assigned in PitchBook’s data; in our broader
PitchBook sample, fewer than 2\% have SIC codes assigned. This low coverage likely reflects that
SIC assignments in commercial databases are concentrated among more established firms.
Batch.
For YC startups, the cohort identifier (e.g., W20, S21, W24) indicating when the startup
was accepted into Y Combinator. Batches occur twice per year (Winter and Summer) plus occasional
Fall batches.
First Deal Quarter.
For PitchBook startups, the year-quarter of the company’s first recorded
financing deal (e.g., 2020-Q1, 2023-Q3). This serves as a timing variable similar to YC batch.
3.4
Research Design
Our research design is descriptive, documenting if AI startups differ from non-AI startups in how
they are organized, a question of significant scholarly and practitioner debate. Our approach follows
in the spirit of recent work establishing stylized facts about who starts successful companies (Azoulay
et al., 2020) and how AI tools are actually used across the workforce (Humlum and Vestergaard,
2025). Here, we provide the first systematic evidence on how firms built around AI are organized.
That said, two confounds are crucial to address to allow for meaningful comparisons between
AI and non-AI startups. First, there is no discrete “arrival” of AI—the term dates to the 1950s.
But since the launch of ChatGPT in late 2022, interest in AI tools and startups has exploded. As
a result, there is the potential for strong cohort confounds: any observed difference between AI and
non-AI firms might reflect when startups were founded rather than their AI status. A startup from
the Winter 24 YC batch is both more likely to be AI and younger when we measure it in January
2025. We address this by including fixed effects for timing. In the case of our YC sample, we use the
startup’s batch, effectively controlling for its quarter of entry. In PitchBook, we don’t have batch,
and instead use first deal quarter.
13
Second, AI startups are likely to concentrate in different industries, so observed organizational
differences could reflect industry differences and market structure rather than AI per se. Enterprise
B2B software firms and consumer hardware firms naturally require different team structures regard-
less of whether they use AI. Yet, building on Chandler (1977), if building around AI causes firms to
operate in different industries, and those industries then impact the nature of the organization, we
could view this selection as an effect of AI. The railroad and telegraph similarly enabled firms to
enter new industries and serve national markets, which in turn shifted their organizational structure.
Here, we take the more conservative approach and focus on within-industry comparisons to better
isolate organizational differences from industry differences.
Our baseline specification for the YC sample is:
Yi = β · 1[AIi] + γbg + εi
(1)
where Yi is a firm-level measure for startup i, 1[AIi] indicates whether the startup is classified as
AI, and γbg are batch × industry fixed effects. For the PitchBook sample, we replace the binary AI
indicator with the continuous AI probability score and replace batch with first deal quarter, using
first deal quarter × industry fixed effects.
By including these interacted fixed effects, we compare AI and non-AI startups operating in
the same industry that entered in the same cohort, ensuring our AI coefficient captures within-
industry-cohort variation rather than cross-industry differences or simple cohort timing confounds.
We report heteroskedasticity-robust (HC1) standard errors throughout.
4
Results
4.1
Summary Statistics and Growth of AI Startups
We begin by documenting basic patterns in our data.
Table 1 presents summary statistics for
our two samples. Even before conditioning on batch or industry, we observe raw differences. AI
startups are substantially smaller than non-AI startups: the mean AI startup in our YC sample
has approximately 13 Revelio-matched employees compared to 42 for non-AI startups (based on
January 2025 snapshots)—a difference of nearly 70\%. AI startups also have flatter hierarchies (3.8
vs. 5.0 seniority levels), fewer managers (17\% vs. 24\%), more engineers (45\% vs. 36\%), and fewer
entry-level workers (19\% vs. 30\%). Despite their smaller size, AI startups have comparable total
14
funding and valuations to non-AI firms, implying that AI startups are substantially more capital-
efficient on a per-employee basis. That said, caution must be taken when comparing these means,
as they fail to control for cohort or industry composition differences. The summary statistics for the
PitchBook sample (Panel B) show similar patterns, though more muted—likely because PitchBook
includes a much broader sample of firms across more industries. This suggests that within-industry
comparisons may be necessary to reveal differences in the broader sample.
Indeed, when we turn to Figure 1 we can see the dramatic rise of AI startups in both samples.
The share of YC startups self-identifying as building around AI has increased substantially, rising
from approximately 15\% in early 2020 batches to over 50\% in 2024 batches. This pattern is mirrored
in our broader PitchBook sample, where the share of early-stage startups we classify as AI has grown
from roughly 8\% in 2020 to over 20\% by 2024. This also underscores the need to compare within
cohorts, as AI startups in our data are from much more recent vintages.
Figure 2 shows the share of AI and non-AI startups across YC sub-verticals (YC’s own finer-
grained and startup-focused industry classification), ordered by relative AI concentration. AI star-
tups cluster in sub-verticals like legal services and developer tools, where applying AI is thought to
have especially large impacts on worker productivity (Eloundou et al., 2024a). Non-AI startups are
more prevalent in hardware-intensive sub-verticals like aviation and medical devices.
4.2
How Are AI Startups Organized?
We first examine whether AI startups have smaller teams. We start with the raw data. Figure 3
shows a histogram of employee counts, confirming that the modal AI startup has fewer than 10
employees, whereas non-AI startups are more evenly distributed across size categories. Given that
there has been an increase in the share of AI startups, these differences could simply be due to
cohort effects. However, in Figure 4 we plot the size of AI and non-AI firms relative to the number
of quarters since they graduated from YC. AI firms are roughly half the size of non-AI ventures
and this gap persists for at least three years. Appendix Figure A3 presents the raw distributions
separately for each YC batch, showing that the distributional difference in Figure 3 between AI and
non-AI startups are evident within each cohort. Finally, these results are not rooted in industry
differences either: Appendix Figure A4 presents kernel density plots of the firm size distribution for
AI and non-AI startups, with batch-industry fixed effects residualized out. The distributions are
visibly shifted, with AI startups concentrated at smaller sizes.
Table 2 reports regression estimates for the relationship between AI status and firm size, com-
15
paring startups within the same cohort and industry. In levels, AI startups in the YC sample have
approximately 8 fewer employees than non-AI startups, compared to a non-AI mean of 34 employees.
In the PitchBook sample, AI startups have approximately 6 fewer employees off a non-AI mean of
18 employees. The log specifications, which examine the intensive margin conditional on having at
least one employee, show that AI startups are roughly 25\% smaller in the YC sample (coefficient of
−0.28) and 12\% smaller in the PitchBook sample (coefficient of −0.13).6 Results are robust across
specifications: levels, logs, and Poisson models all show significant negative effects of AI status on
firm size. Appendix Table A4 shows these patterns hold when we use alternative measures of firm
size sourced from PitchBook and YC rather than Revelio.
These teams are not just smaller but employ a different mix of functional expertise. Table 3
examines the composition of startup workforces across seven functional categories.
AI startups
employ a significantly higher share of technical workers relative to business and administrative
staff. The coefficient on AI status when predicting engineering share is positive and significant in
both samples (β = 0.045, p < 0.01 in YC; β = 0.183, p < 0.01 in PitchBook), indicating that AI
startups allocate roughly 5 percentage points more of their workforce to engineering roles in YC and
18 percentage points more in PitchBook. Conversely, AI startups have significantly lower shares
of employees across business functions: sales (β = −0.021, p < 0.05), operations (β = −0.009,
p < 0.05), finance (β = −0.017, p < 0.01), and administrative roles (β = −0.011, p < 0.05).
The coefficient on marketing share is negative but not statistically significant in the YC sample
(β = −0.012). The statistically significant declines span both external-facing functions (sales) and
internal process functions (finance, operations, administration) in YC, and extend to all business
functions including marketing in the broader PitchBook sample.
This suggests that the talent
needed to scale firms may shift when they are built around AI (Bresnahan, Brynjolfsson, and Hitt,
2002; DeSantola, Gulati, and Zhelyazkov, 2022; Lee and Kim, 2024). These patterns are even more
pronounced in the broader PitchBook sample, where all functional coefficients are significant at
the 1\% level. We interpret these cautiously, however, because AI status in this sample is inferred
from a text-based prediction, which may mechanically inflate compositional differences. The picture
that emerges is that AI startups are technical-heavy and business-function-light—consistent with
shifting the firm’s production function toward the technical talent needed to build and maintain AI
6For small coefficients, eβ −1 ≈β, so log coefficients can be read directly as approximate percentage changes. For
larger coefficients, the exponential transformation is more accurate. Throughout the paper, we convert log
coefficients to percentage changes using the transformation eβ −1, which provides an exact interpretation for
discrete changes in AI status.
16
capabilities. That said, it is important to remember that these firms are smaller, thus while there
are relatively more engineers, there are fewer employees in each role across the board for AI-native
firms.
We next turn to workforce seniority. Table 4 examines the seniority composition of AI versus
non-AI startups. AI startups have 5 percentage points more senior employees (β = 0.052, p < 0.01)
and 4 percentage points fewer entry-level employees (β = −0.040, p < 0.01). The broader PitchBook
sample shows the same pattern, with PB AI probability associated with a 6 percentage point lower
entry-level share (β = −0.060, p < 0.01) and a 6 percentage point higher senior share (β = 0.056,
p < 0.01).
This is consistent with Brynjolfsson, Chandar, and Chen (2025), who document a
13\% relative decline in early-career employment in AI-exposed occupations. At the task level, a
growing body of evidence suggests that AI often disproportionately benefits less experienced workers
(Brynjolfsson, Li, and Raymond, 2023; Dell’Acqua et al., 2023; Noy and Zhang, 2023; Csaszar,
Ketkar, and Kim, 2024). One might therefore expect AI-native firms to employ more junior talent.
We find the opposite: AI-native firms employ fewer entry-level workers and more senior workers.
This contrast underscores the importance of studying AI at the organizational level. Who benefits
from AI depends not only on task-level productivity gains, but also on where AI is applied in the
firm’s production process (Gans and Goldfarb, 2026; Agrawal, Gans, and Goldfarb, 2024; Benzell
and Myers, 2026; Otis et al., 2025; Kang and Kim, 2025).
Given that the teams are smaller, more senior, and more technical, we also examine whether
hierarchy differs. Figure 5 shows the distribution of hierarchy depth by AI status. AI startups
have significantly flatter organizational structures: the median AI startup has 3 distinct seniority
levels compared to 4 for non-AI startups. Table 5 confirms this pattern in regression estimates,
with AI startups having roughly 0.5 fewer hierarchy levels on average (β = −0.482, p < 0.01). In
Appendix Table A5 we show these findings are robust to controlling for firm size—the flattening
is not simply a mechanical consequence of being smaller (see also Appendix Figure A5 for the
residualized distribution).
Consistent with less need for coordination in AI startups, these firms also have a lower share
of employees in management occupations. Using O*NET occupation codes to identify managers
(codes beginning with “11-”), we find that AI startups have approximately 4 percentage points fewer
managers as a share of total employment (β = −0.037, p < 0.01). These results are reminiscent
of the flattening of firms due to information technology (Brynjolfsson et al., 1994; Bresnahan,
Brynjolfsson, and Hitt, 2002; Rajan and Wulf, 2006; Bloom et al., 2014), suggesting that when
17
productive capability moves from people to the product, the need to coordinate and supervise
internal knowledge work may diminish.
4.3
Do AI Startups Achieve More with Less?
The preceding findings raise a natural concern: perhaps AI startups are smaller and flatter simply
because they are lower quality or have less potential. AI may lower entry barriers, allowing anyone
to enter, even with a weak idea (Reimers and Waldfogel, 2026). If so, the organizational differences
we document would reflect the limitations of building around AI rather than evidence that AI is
changing the nature of the firm.
We begin by examining total funding and valuation in levels. Table 6 (Panel A) reports regression
estimates. In the YC sample, AI startups raise approximately 15\% less total funding (β = −0.165,
p < 0.05) and show no statistically significant difference in valuations. In the broader PitchBook
sample, the pattern reverses: AI startups raise approximately 15\% more funding (β = 0.136, p <
0.10) and achieve 58\% higher valuations (β = 0.458, p < 0.01). This divergence likely reflects
YC’s selection: virtually all YC startups are technology firms that have met a high admission bar,
compressing differences in funding and valuation relative to the broader PitchBook universe.
The more informative comparison is value per employee. In Figure 6, the raw distributions of
valuation per employee show that AI startups are shifted rightward, achieving higher valuations
relative to their headcount. Appendix Figure A6 residualizes out batch fixed effects and shows a
similar pattern.7
Table 6 (Panel B) quantifies these patterns, focusing on the subset of startups with available
funding and valuation data in PitchBook.8 In the YC sample, AI startups raise approximately 12\%
more funding per employee (β = 0.109, p < 0.10) and result in 30\% higher valuation per employee
(β = 0.260, p < 0.05). The PitchBook sample confirms these patterns more strongly: AI startups
raise 33\% more per employee (β = 0.285, p < 0.01) and achieve 76\% higher valuation per employee
(β = 0.564, p < 0.01). These efficiency results are consistent across both samples and robust to
specification.
Notably, the valuation-per-employee effects are substantially larger than the funding-per-employee
7Valuations reflect investors’ expectations of future cash flows rather than realized value. We use valuation per
employee as a measure of capital-market-implied productivity, not as a direct measure of output.
8Funding-per-employee results are robust to treating missing funding values as zeros, which has a defensible
interpretation (no recorded deal). For valuation, missing values reflect “no priced round on record” rather than
literal zero value (no firm in our sample has a PitchBook-recorded valuation of exactly zero), so we restrict the
valuation sample to firms with a non-missing valuation rather than imputing zeros. Poisson estimates yield similar
conclusions for funding outcomes.
18
effects in both samples. If AI startups were simply raising more capital, we would expect funding
and valuation to increase proportionally. Instead, the disproportionate increase in valuation rela-
tive to funding is consistent with investors expecting higher future productivity from AI startups
conditional on the capital and labor inputs.
4.4
Where Do AI Startups Locate and Who Works for Them?
A large literature demonstrates that as the nature of organizations changes, so too do the loca-
tions where firms concentrate and the workers they employ (Botelho, Gulati, and Sorenson, 2024).
New technologies can have disproportionate impacts on particular groups—shaping employment
opportunities and potentially long-term career prospects (Feigenbaum and Gross, 2022). We briefly
examine these patterns for AI startups, reporting figures and tables in the Appendix.
Appendix Table A8 shows that AI startups are more geographically concentrated than non-AI
startups. AI startups are 20 percentage points more likely to be headquartered in San Francisco
(p < 0.001), a substantial increase relative to the baseline rate of just under a quarter among non-
AI startups. This concentration likely reflects the importance of access to specialized AI talent
and proximity to leading AI research labs, consistent with longstanding findings on the distinctive
advantages of Silicon Valley as a technology hub (Guzman and Stern, 2015; Saxenian, 1994). These
geographic patterns hold in the broader PitchBook sample as well (Appendix Table A11).
We also find demographic and educational differences between AI and non-AI startup workforces
(Appendix Tables A6–A10).
These demographic patterns are consistent with the expert-dense
organizational form we document: the shift toward technical, senior workforces draws from labor
pools that remain less diverse. In the YC sample, AI startup employees are about 5 percentage
points more likely to hold graduate degrees (Appendix Table A7). In the broader PitchBook sample,
AI probability is associated with a 5.9 percentage point lower female share, an 8.5 percentage point
higher Asian/Pacific Islander share, a 17.6 percentage point higher graduate-degree share, and
higher university prestige (Appendix Tables A9 and A10).9
These geographic and demographic patterns raise questions about whether the benefits of AI-
enabled organizational forms will accrue evenly, given persistent disparities in entrepreneurial finance
by gender and race (Ewens, 2022; Marx, Wang, and Yimfor, 2025). Evidence suggests that women
9Revelio Labs computes prestige scores using an expectation-maximization model that alternates between user and
institution prestige updates, with world university rankings as priors. Prestige values are continuous between −1
and 1, weighted by seniority and tenure. See https://www.data-dictionary.reveliolabs.com/methodology.html
for full methodology.
19
adopt generative AI at lower rates than men across multiple contexts, even after controlling for
occupation, education, and technical skill (Otis et al., 2026). If AI tools accelerate learning for
those who use them, differential adoption rates may translate into widening performance gaps—
both for individual workers within firms and for the entrepreneurs who found them.
5
Exploring AI Product and AI Process
The results so far establish that AI-native firms are smaller, flatter, more technically concentrated,
and more capital-efficient than non-AI startups in the same industry and cohort. But these patterns
alone do not tell us why AI startups are organized differently. As discussed in Section 2, we note
that two primary channels could account for these patterns: workers using AI tools to do their
jobs more productively (the process channel), and AI capabilities embedded directly into the firm’s
product (the product channel). In this section we measure each channel separately and explore how
each accounts for the organizational differences we documented above.
5.1
Measuring AI in a Firm’s Product and Processes
To quantitatively measure the product channel, we begin by classifying how founders describe the
role of AI in their businesses. To do so, we built an inductive classification of the descriptions
founders post of their startup on YC’s startup directory platform. To generate the classification,
we first drew a random sample of 40 AI-tagged startup descriptions and read them qualitatively,
identifying recurring patterns in how AI enters the product. We formalized those patterns into
four categories and then used an LLM to apply the schema to the 990 AI-tagged YC startups (see
Appendix for the full methodology and prompt). Figure 7 presents the resulting distribution of
AI-tagged startups, and Table 7 provides concrete examples from each category. For the regression-
based AI Product measure we use below, we re-apply the same schema to the full YC sample so AI
and non-AI firms are coded symmetrically.
The classification captures different ways that startups embed AI into the product. The largest
category, autonomous workflow products (43\%), includes startups whose products use AI to perform
tasks or roles that previously required human labor. Artisan (W24) builds AI sales representatives
that discover leads, research prospects, and write personalized emails. Aurelian (S22) deploys AI
phone operators that triage 911 calls. These products place AI directly in the workflow, allowing the
product itself to perform work that would otherwise require salespeople or operators. The startup
20
can offer the service without building a large internal organization to deliver the underlying work
itself.
The second category, expert augmentation products (24\%), includes startups whose products
use AI to extend a professional’s existing workflow. Tower (W24) helps lawyers stay organized and
close deals faster. Medium Biosciences (S21) uses AI to design therapeutic proteins, augmenting
biologists in drug development. Here, AI does not replace the expert role outright. Instead, the
product embeds analytic, generative, or search capabilities that expand what expert users can do.
This also moves productive capability into the product: rather than hiring large internal teams of
lawyers or scientists to provide expert services directly, the startup builds a product that gives users
some of those capabilities themselves.
These first two categories are clear examples of startups embedding AI into their products.
Whether the startup is augmenting or automating work, in both cases the company is building a
product around AI. Taken together, we find that about 67\% of YC startups tagged as AI-native are
clearly embedding AI directly into the products they sell.
The third category, AI infrastructure products (15\%), includes firms whose products help others
build, deploy, or use AI. Exa (S21) offers an AI-powered search engine using embeddings and neural
retrieval, and Roboflow (S20) provides a computer vision platform for deploying object detection
models. These firms embed AI into tools that become part of other companies’ production systems.
The product does not always automate an end-user workflow directly, as in autonomous workflow
products, or augment a specific professional role, as in expert augmentation products. Instead, some
packages model capabilities, data pipelines, or deployment tools into infrastructure that other firms
can build on. Thus, some of these firms clearly embed AI in the products they sell while others sell
infrastructure that complements the use of AI (e.g., monitoring and analytics tools).
Finally, the remaining 18\% (“Other”) includes startups where AI’s role is present but harder to
classify, including games, robotics platforms, or firms whose descriptions do not clearly specify how
AI is embedded in the product.
To create a proxy for the process channel, we operationalize internal AI tool use directly from
firms’ job posting text. For each firm in our sample we collect all Revelio Labs postings dated
through January 2025—the same window over which we measure employment—and drop scrapes
flagged as staffing-firm aggregator listings, instead focusing on postings directly listed by the hiring
firm. We then code a firm-level binary indicator, AI Process, equal to one if any of the firm’s
non-staffing postings names a specific Gen AI assistance tool from a dictionary of 33 AI tools. The
21
dictionary includes general-purpose AI assistants (e.g., ChatGPT, Claude), coding agents (e.g.,
GitHub Copilot, Cursor), creative tools (e.g., Midjourney, Adobe Firefly) and workflow products
(e.g., Granola, Mem.ai). The dictionary was assembled by reading the corpus of job postings and
noting listed AI tools along with tools included in the Stack Overflow 2024 Developer Survey of
AI-tools section. The full dictionary and construction details are listed in Appendix A.3.
Two example postings from AI-tagged YC firms illustrate why focusing on tools, rather than
generic mentions of “AI,” is more likely to capture process-focused uses. The first, from Overlap
(S24), a multimodal-AI firm hiring a Product Engineer, is saturated with AI language—“multimodal
AI agents,” “Video Marketing AI,” “AI-powered products,” “AI-driven media”—but every reference
describes the AI features that will be built into the product. No specific AI tool, model, or assistant
from our dictionary is named, so AI Process equals zero. The second, from Fintool (W23) hiring a
Full-Stack Engineer, contains a clean process-channel signal: “We are heavy users of V0, Cursor and
other AI tools.” That sentence describes the firm’s workers using specific named AI tools to ship
product features, and it matches the V0 and Cursor patterns in our dictionary, so AI Process equals
one. The channel regressions use the 1,684 firms with at least one non-staffing posting through
January 2025; the log, hierarchy, manager-share, and base-rate analyses further require at least one
Revelio-matched employee, leaving 1,666 firms (467 AI-tagged and 1,199 non-AI startups).
Within this sample of 1,666 firms with at least one job posting, the product channel is present
at 69.8\% of AI-tagged firms versus 16.1\% of non-AI firms, while the process channel is present at
13.5\% of AI-tagged firms versus 5.2\% of non-AI firms—a roughly 2.6× relative gap on the process
side, though with substantially smaller absolute levels than the product channel. These lower levels
likely reflect both the nascency of AI assistance tooling over our sample window and th
\end{document}
